% from main.tex
%%%%%%%%%%  scan idx 164  |  folio 151  %%%%%%%%%%
\emph{Preuve.} On définit la flèche en résolvant \(F\) injectivement et en utilisant la flèche canonique de \(\Kpx(S)\)
\[
  E\longrightarrow\Hom^\bullet(\Hom^\bullet(E,F),F).
\]
Pour voir que la flèche en question est un isomorphisme si \(E\) est parfait et \(F=A\), on se ramène, par localisation et dévissage, à \(E=A\), auquel cas c'est évident.

\textbf{Définition 7.3.} \emph{Soit \(E\in\Ob\Dpx(S)\). Nous poserons \(\RHom(E,A)=E^\vee\), et nous dirons que \(E^\vee\) est le dual de \(E\).}

D'après (7.1) et (7.2), si \(E\in\Ob\Dpx(S)\) est parfait (resp. parfait et de tor-dimension finie), \(E^\vee\) l'est également, et la flèche canonique \(E\to E^{\vee\vee}\) est un isomorphisme. En d'autres termes, le foncteur \(\RHom(\,\cdot\,,A)\) définit une anti-auto-équivalence de la catégorie \(\Dpx(S)_{\mathrm{parf}}\) (resp. \(\Dpx(S)_{\parf}\) (5.8.4.b))) avec elle-même.

\textbf{Lemme 7.4} (isomorphisme de CARTAN) (cf. ([H] II 5.15)). Pour \(E\in\Ob\Dpx(S)\), \(F\in\Ob\Dpx^-(S)\), \(G\in\Ob\Dpx^+(S)\), il existe des isomorphismes fonctoriels canoniques :
\[
\begin{aligned}
\RHom(E\otL F,G)&\xrightarrow{\sim}\RHom(E,\RHom(F,G)),\\
\RHom(E\otL F,G)&\xrightarrow{\sim}\RHom(E,\RHom(F,G)),\\
\Hom(E\otL F,G)&\xrightarrow{\sim}\Hom(E,\RHom(F,G)).
\end{aligned}
\]

